% =====================================================================
% Methodology , revision draft, 2026-08-12.
%
% Every number traces to a committed artifact named in the comment above it,
% or to a \FN* macro defined in sec_frozen_numbers.tex. Do not restate a number
% from memory: cite the file. Experimental record CLOSED 2026-08-13.
% =====================================================================

\section{Methodology}

We conduct a controlled ablation that isolates the two design choices most
directly shaping a spiking intrusion detector: the \emph{neuron model} and the
\emph{spike-encoding strategy}. Every other pipeline component is held fixed.

\subsection{Scope}

This study is a controlled intra-SNN design-space evaluation. The nine neuron
families are \texttt{snnTorch} implementations; the three spike encoders are
our own and are released with the artifact. Every model evaluated is a spiking network; the
axes of variation are the neuron family and the spike encoding. We deliberately
do not include non-spiking baselines, decision-tree ensembles, linear models,
or an architecture-matched dense MLP. Such a comparison answers a different
question, namely whether spiking networks are competitive with conventional
tabular classifiers, and answering it credibly requires its own tuning protocol
and its own fairness argument. Our question is narrower and, we argue, logically
prior: \emph{given} that a spiking detector is to be deployed, which neuron and
encoding should it use? We therefore make no claim of superiority over
conventional machine learning, and none of our efficiency results should be read
as an SNN-versus-DNN comparison. The efficiency analysis (Section~\ref{sec:eff})
stays inside the spiking design space for the same reason.

\subsection{Pre-registration}

Because this revision re-evaluates a configuration we previously identified as
best, the ranking rules were frozen in a pre-registration document before any
statistic in the revised analysis was computed. It fixes macro-F1 as the primary
quality metric, the (dataset $\times$ seed) pair as the experimental unit, mean
rank as the primary ranking statistic, the statistical tests, and, critically, a mapping from outcome to permitted manuscript wording, so that the strength
of the claim follows the evidence rather than the reverse. It also forbids
several specific practices: selecting datasets after seeing the ranking,
choosing operating thresholds on test data, and reporting a pooled
cross-dataset standard deviation. The document is included with the artifact.

\subsection{Two protocol generations}

We report two evaluation protocols and never mix them within a table.

\textbf{The screening protocol} is the one used for the full sweep. Its role is the full
$9\times3=\FNScrVariants$ configuration \emph{screening} sweep, and we preserve
it unchanged so that the previously published numbers remain reproducible. We
use \emph{screening} and \emph{strict confirmation} consistently for the two
generations, because their populations differ and a reader must be able to tell
at a glance which one a statistic belongs to.

\textbf{v2} is leakage-resistant and its role is \emph{confirmation}. It differs
from v1 in four respects, each addressing a specific methodological weakness:

\begin{enumerate}
  \item \textbf{Nothing is fitted on test.} Categorical vocabularies, the
        $\log(1+x)$ transform and min--max parameters are fitted on training
        rows only, with an explicit \texttt{<unk>} column per categorical field
        for categories first seen at evaluation time. v1 one-hot encoded against
        the union of all splits. We report the size of that earlier leak rather
        than merely asserting it was small: it affected $0$ NSL-KDD and $2$
        KDDCup99 test records. It was therefore empirically negligible on these
        datasets, but the revised protocol nonetheless fits every learned
        transform on training data alone.
  \item \textbf{The partition is frozen independently of the model seed.} In v1
        the CIC-IDS2017 and CTU-13 partitions were drawn using the model seed,
        so seeds $42$--$46$ varied the partition and the initialisation
        together and their spread conflated two sources of variance. In v2 the
        partition is a property of the protocol; the seed only initialises
        weights and orders batches.
  \item \textbf{Group structure is respected} where the benchmark has one:
        capture file for CIC-IDS2017, scenario for CTU-13.
  \item \textbf{Every split is audited before use.} A split that declares a
        group-disjointness it does not satisfy raises an exception and the
        training job refuses to start.
\end{enumerate}

% results/analysis/split_audit.csv
Table~\ref{tab:audit} reports the audit. Alongside disjointness it records exact
duplicate rows shared between training and test. These are reported, never
asserted to be zero: $57.53\%$ of the official KDDCup99 test set occurs verbatim
in training, which is the documented reason NSL-KDD exists. We accompany every
KDDCup99 result with a duplicate-free sensitivity figure computed on the
non-duplicated remainder.

\subsection{Datasets and partitions}

% results/manifests/*.json
NSL-KDD and KDDCup99 use their official train/test splits. For CIC-IDS2017 we
adopt a \emph{within-capture order-disjoint} partition: within each of the eight
capture files the first $70\%$ of rows form the training pool, the next $10\%$
validation, and the final $20\%$ test, without shuffling, so that flows from one
attack burst are not scattered across the evaluation boundary. We describe this
as order-disjoint rather than chronological because the public redistribution we
use does not retain the original \texttt{Timestamp} field; row position is
CICFlowMeter's write order and serves as an ordering proxy. The day-level
grouping, which carries the genuine temporal claim, is exact because it derives
from the capture file names. We additionally evaluate a \emph{day-disjoint}
protocol, train on Monday--Wednesday, test on Thursday--Friday, as a
binary normal-versus-attack transfer test. It is binary by necessity: under the
five-class taxonomy, Probe and U2R occur only in the later days and R2L only in
the earlier ones, so three of five classes would be untrainable or untestable.

The two confirmation folds were selected before the confirmation registration
was frozen and before any confirmation result existed, using per-fold difficulty
observed in a prior, separate experiment. Fold~1 is the hardest of the four
($\FNCtuIdxFB\%$ macro-F1 under the four-fold performance estimate) and fold~0
is near the ceiling ($\FNCtuIdxFA\%$), so the confirmation spans both regimes
rather than a middling pair. This is a difficulty-informed choice, not a blind
one, and we state it as such: the registration governs the confirmation
hypotheses, metric, ranking rule and permitted wording, not the fold selection,
which preceded it.

For CTU-13 we use a frozen four-fold \emph{scenario-disjoint} partition in which
whole scenarios are indivisible, and we recompute the per-source-host behavioural
statistics \emph{causally}: within a scenario, flows are processed in timestamp
order and each flow's host features derive from strictly earlier flows of the
same source only. The screening protocol computed those aggregates over an entire
scenario and then split flows at random, so a test flow's features were computed
partly from test flows, a feature-level leak independent of how rows are
partitioned. We emphasise that the scenario protocol is \textbf{not}
host-disjoint: $76.2\%$ of the test source hosts in fold~1 also appear in
training, and we report it as scenario-disjoint only.

\subsection{Two-axis confirmation design}

A complete v2 re-sweep of all 27 configurations costs approximately $198$
GPU-hours on our hardware and was outside the revision budget. Rather than
re-running an arbitrary subset, we confirm the selected configuration along both
design axes that pass through it:

\begin{itemize}
  \item \textbf{Axis A (encoding).} \texttt{LeakyParallel} held fixed, all three
        encodings varied.
  \item \textbf{Axis B (neuron).} Latency encoding held fixed, all nine neuron
        families varied.
\end{itemize}

Their intersection is \texttt{LeakyParallel/latency}. Axis~B is a confirmation
that the neuron choice survives a stricter protocol; it is not a second
design-space ranking and we do not present it as one.

\subsection{Spike encodings and their amplitude gates}
\label{sec:enc}

Each feature vector $x\in[0,1]^d$ is encoded over $T$ timesteps. \textbf{Rate}
coding draws $S_{t,i}\sim\mathrm{Bernoulli}(x_i)$ independently at each
timestep. \textbf{Latency} coding assigns each feature above an amplitude gate
$\theta_\mathrm{lat}=\FNLatThresh$ a single spike at
$t_i=\lfloor(1-x_i)(T-1)\rfloor$. \textbf{Delta} coding is applied to a
synthesised step signal in which $x$ is held for the first $\lfloor T/2\rfloor$
timesteps and zero thereafter, with negative transitions discarded. The
encoders are our own; only the neuron families come from \texttt{snnTorch}, and
all three are released with the artifact.

Two of the three apply an amplitude gate, at different thresholds, which could
in principle confound the encoding comparison with a thresholding effect. We
measured the affected feature mass rather than assuming it small. The interval
$(0,\theta_\mathrm{lat}]$ is \emph{empty} on both KDD-family protocols
($\FNBandNsl$ and $\FNBandKdd$ features per sample) and contains
$\FNBandCic$ of $\FNDimCic$ features on CIC-IDS2017 and $\FNBandCtuA$ of
$\FNDimCtuA$ on CTU-13. Latency's gate is therefore inert: setting
$\theta_\mathrm{lat}=0$ leaves the measured input-spike count unchanged to four
decimals ($\FNNslLatInSpk$ on NSL-KDD in both cases). The interval separating
the two gates, $(\theta_\mathrm{lat},\theta_\delta]$, contains approximately one
of the $\FNNslLatInSpk$ active NSL-KDD features, and
Section~\ref{sec:e1} reports a gate-matched sensitivity analysis over that
residual difference.

\subsection{Metric definition and a post-audit correction}

Macro-F1 is computed over the \emph{fixed} label universe of the task, so a
class absent from a test partition contributes zero rather than being dropped
from the average. This matters on CIC-IDS2017, whose order-disjoint test
interval contains no U2R samples: the metric there is capped at
\FNCicCeiling{}, and its absolute value is not comparable to a five-class score
on a dataset where all five classes occur.

An audit after the runs completed found that our metric implementation had
omitted an explicit label set, so the averaging denominator was the union of
labels appearing in the references or the predictions. On the affected protocol
that made the denominator depend on the model: a configuration that never
predicts the unsupported class was averaged over four classes while one emitting
a single false positive for it was averaged over five. We corrected the
implementation and recomputed every affected figure from the persisted
per-sample predictions, without retraining. The correction leaves every ranking,
mean rank and hypothesis test in this paper unchanged; it alters one absolute
value. Alongside the pre-registered primary metric we report macro-F1 over the
classes represented in test as a descriptive quantity, which is never used for
selection or testing. The full record is in the artifact's amendment log.

\subsection{Statistical protocol}

% results/analysis/v1/{friedman,focused_tests,encoding_tests}.json
Each (dataset, seed) pair is one paired block, giving $4\times5=20$ blocks.
Within each block all configurations are ranked by macro-F1. The primary
statistic is the \emph{mean rank} across blocks rather than the mean of raw
macro-F1 across datasets, because a raw cross-dataset mean is dominated by how
different the datasets are from one another rather than by how the
configurations differ. For the same reason, dispersion is always reported within
a dataset as mean $\pm$ standard deviation over seeds; we do not report pooled
cross-dataset standard deviations.

We test for an overall difference with a Friedman test across configurations
and, where it rejects, apply a Nemenyi post-hoc procedure and report the
critical difference. Two focused comparisons were specified in advance:
\texttt{LeakyParallel/latency} against \texttt{LeakyParallel/rate}, its closest
competitor, and against the best non-\texttt{LeakyParallel} configuration. These
use a paired Wilcoxon signed-rank test over the blocks with Holm correction
across the family, and are accompanied by a hierarchical bootstrap interval that
resamples datasets first and seeds within datasets. Where the block-level test
and the dataset-level interval disagree, we report both and say so: that pattern
indicates an effect consistent in direction but variable in magnitude across
datasets, which is a different and weaker claim than uniform superiority.

\subsection{Anytime readout}

% results/analysis/early_readout/
Because latency coding presents high-magnitude features earlier in the window,
terminal accuracy alone may understate its value. For every saved model we
therefore evaluate the readout after each intermediate timestep,
$\hat{y}(t)=\arg\max_c\sum_{\tau\le t}S_c(\tau)$, and record macro-F1, MCC,
detection rate and cumulative spike activity as functions of $t$. We define
$T_{95}$ and $T_{99}$ as the first timestep at which macro-F1 reaches $95\%$ and
$99\%$ of its terminal value. Both thresholds were fixed in advance. This
analysis requires no additional training: the per-timestep cumulative output
counts are persisted for every run.

\subsection{Efficiency accounting}
\label{sec:eff}

% results/analysis/efficiency/summary.csv
We decompose spike activity into input, hidden and output components rather than
reporting a single aggregate, and weight each spike by the fanout it drives:
$\mathrm{SOP}=\sum_{\ell,t}s_\ell(t)\cdot\mathrm{fanout}_\ell$. The decomposition
matters because an undifferentiated spike count can exceed the input
dimensionality, leaving the reader unable to tell what it contains; the fanout
weighting matters because a spike entering a 128-wide projection is not the same
event as one entering a 5-wide readout. Gated cells (\texttt{SLSTM},
\texttt{SConv2dLSTM}) are charged four times the input fanout to account for
their gate projections. Where we convert operations to energy we report
SOPs as a hardware-independent operation count. We deliberately do not convert
to joules: the conversion constant is platform-specific and quoting one would
imply a hardware measurement we did not make. Output-layer spikes are charged
zero fanout, because the readout drives no downstream projection, so the count
comprises input-to-hidden and hidden-to-output operations only.

\subsection{Experimental setup}

All experiments use PyTorch~2.12 with \texttt{snnTorch}~0.9.4 on a
single NVIDIA H200 GPU. The network is a two-layer spiking multilayer perceptron
with hidden width $H=128$, unrolled for $T=25$ timesteps, membrane decay
$\beta=0.85$, synaptic decay $\alpha=0.90$ for the conductance-based families,
and the arctangent surrogate gradient. Training uses Adam at $10^{-3}$ for
$10$ epochs with batch size $128$, minimising cross-entropy on the accumulated
output-layer spike counts. No per-variant hyperparameter tuning is performed.

The two gated sequence neurons use a reduced firing threshold of $0.1$; at the
library default of $1.0$ they emit no spikes on sparse tabular input and
collapse to the majority class. This means the sweep is not a pure neuron-only
ablation, and we state that limitation rather than presenting the comparison as
threshold-matched (Section~\ref{sec:threats}).

Each configuration is repeated over five seeds $\{42,\dots,46\}$. We seed NumPy,
PyTorch and CUDA and enable deterministic cuDNN kernels. Every result record
carries the identity of the frozen split that produced it, including a hash of
the split manifest, so any reported number can be traced to its exact partition.
